% F085｜真实数学图；代表函数不构成抽象解的唯一性证明。
\begin{center}
\begin{tikzpicture}[x=2.3cm,y=2.3cm,>=stealth]
\draw[->] (-.1,0)--(2.1,0) node[right,font=\small] {$u$};
\draw[->] (0,-.1)--(0,1.8) node[above,font=\small] {$v$};
\draw[CExplanation,thick] (1,0) arc(0:90:1);
\draw[black!40] (1,0)--(1,1.7);
\draw[black!40] (0,1)--(1.9,1);
\draw[CExplanation,very thick] (0,0)--(1.7,1.7);
\fill[CExplanation] (.7071068,.7071068) circle(.022) node[above left,font=\small] {$P$};
\draw[dashed] (.7071068,0)--(.7071068,.7071068)--(0,.7071068);
\draw[orange!80!black,very thick] (1,0)--(1,1);
\draw[green!45!black,very thick] (0,1)--(1,1);
\draw (.25,0) arc(0:45:.25);
\node[font=\small] at (.38,.13) {$\theta$};
\node[font=\small,anchor=west,text=orange!80!black] at (1.07,.5) {$\tan\theta$};
\node[font=\small,above,text=green!45!black] at (.5,1.03) {$\cot\theta$};
\node[font=\small,anchor=west,align=left] at (2.6,1.2) {$P=(\cos\theta,\sin\theta)$\\[7pt]
竖直切线 $u=1$：\\[3pt]$\tan\theta=\sin\theta/\cos\theta$\\[7pt]
水平切线 $v=1$：\\[3pt]$\cot\theta=\cos\theta/\sin\theta$};
\end{tikzpicture}
\par\small 取 $\theta=x+y$。同一条射线的纵横比与横纵比互为倒数；相应分母为 $0$ 时，比值无定义。图示 $\theta=\pi/4$。
\end{center}
